\paragraph{Learning to choose a complete drawing.}
We tested whether a writer learns the next glyph more effectively when it also has to distinguish the target drawing from seven alternatives. Three further fits used the same parent, minibatch order and 1,920 updates. On 2,048 requests from 1,020 previously unused chapters, the likelihood control recovered four complete modern translations. Two required changing the source, including a classical question rewritten in modern Chinese (Fig.~\ref{fig:edit-contrast}a). The complete-glyph contrast recovered two sentences, both verbatim copies. It increased the aggregate edit rate from 67.90\% to 69.92\% (paired chapter-bootstrap 95\% interval for the increase, 1.75--2.29 percentage points). Of 8,025 independently drawn glyphs from the control, two matched no bitmap in the scoring collection. The corresponding count for the contrastive writer was 6,982 of 8,760. All 2,043 renderable quotations were exact in both conditions.

The contrastive score distinguishes the correct glyph from sampled alternatives. We tested whether correcting the score for their unequal frequencies could recover the quality of free writing. The correction reduced the fresh edit rate to 68.81\%, but recovered neither changed sentence. In a further 755 previously unused chapters the corrected editor and likelihood control each returned seven exact translations on 2,048 requests. Their edit rates were 70.25\% and 69.23\%, respectively (paired increase, 0.85--1.20 percentage points). Local edit counts also gave 69.23\% (Fig.~\ref{fig:edit-contrast}b).

The likelihood control used source order. Supplying correct previous writing, then reversing the source while keeping the same glyph under the cursor, reduced exact drawings from 5,633 to 4,695 of 55,758 positions (paired accuracy loss, 1.49--1.88 percentage points). In 6,847 natural cross-chapter pairs with matched cursor and recent history, the control preferred the correct assignment of the two offered glyphs in 3,963 pairs. It independently drew both target images in 15 pairs. The contrastive writer preferred 3,985 correct assignments but drew both in one pair. Distinguishing two plausible readings gives limited evidence about which drawing the unrestricted writer will produce.
